% TODO proof-reading
\subsubsection{簡単な例}

まず、10個の整数の配列を作成して、それを満たしてみましょう：

\begin{lstlisting}[style=customjava]
	public static void main(String[] args) 
	{
		int a[]=new int[10];
		for (int i=0; i<10; i++)
			a[i]=i;
		dump (a);
	}
\end{lstlisting}

\begin{lstlisting}
  public static void main(java.lang.String[]);
    flags: ACC_PUBLIC, ACC_STATIC
    Code:
      stack=3, locals=3, args_size=1
         0: bipush        10
         2: newarray       int
         4: astore_1      
         5: iconst_0      
         6: istore_2      
         7: iload_2       
         8: bipush        10
        10: if_icmpge     23
        13: aload_1       
        14: iload_2       
        15: iload_2       
        16: iastore       
        17: iinc          2, 1
        20: goto          7
        23: aload_1       
        24: invokestatic  #4     // Method dump:([I)V
        27: return        
\end{lstlisting}

\TT{newarray}命令は、10個の\IT{int}要素の配列オブジェクトを作成します。

配列のサイズは\TT{bipush}で設定され、\ac{TOS}に残されます。

配列の型は\TT{newarray}命令のオペランドで設定されます。

\TT{newarray}の実行後、ヒープに新しく作成された配列への\IT{参照}（またはポインタ）が\ac{TOS}に残されます。

\TT{astore\_1}は、\IT{参照}を\ac{LVA}の1番目のスロットに格納します。

\main関数の第二部分は、\IT{i}を対応する配列要素に格納するループです。

\TT{aload\_1}は配列の\IT{参照}を取得し、スタックに置きます。

\TT{iastore}は、スタックから整数値を取得し、現在\ac{TOS}にある\IT{参照}の配列に格納します。

\main関数の第三部分は、\TT{dump()}関数を呼び出します。

その引数は\TT{aload\_1}（オフセット23）で準備されます。

次に\TT{dump()}関数に進みましょう：

\begin{lstlisting}[style=customjava]
	public static void dump(int a[])
	{
		for (int i=0; i<a.length; i++)
			System.out.println(a[i]);
	}
\end{lstlisting}

\begin{lstlisting}
  public static void dump(int[]);
    flags: ACC_PUBLIC, ACC_STATIC
    Code:
      stack=3, locals=2, args_size=1
         0: iconst_0      
         1: istore_1      
         2: iload_1       
         3: aload_0       
         4: arraylength   
         5: if_icmpge     23
         8: getstatic     #2      // Field java/lang/System.out:Ljava/io/PrintStream;
        11: aload_0       
        12: iload_1       
        13: iaload        
        14: invokevirtual #3      // Method java/io/PrintStream.println:(I)V
        17: iinc          1, 1
        20: goto          2
        23: return        
\end{lstlisting}

入力配列への\TT{参照}は0番目のスロットにあります。

ソースコード中の\TT{a.length}式は\TT{arraylength}命令に変換されます：
これは配列への\IT{参照}を取り、配列サイズを\ac{TOS}に残します。

オフセット13の\TT{iaload}は配列要素の読み込みに使用され、
配列の\IT{参照}がスタックに存在する必要があります（11の\TT{aload\_0}で準備）、
そしてインデックスも必要です（オフセット12の\TT{iload\_1}で準備）。

言うまでもなく、\IT{a}で始まる命令は誤って\IT{配列}命令として理解される可能性があります。

これは正しくありません。
これらの命令はオブジェクトへの\IT{参照}で動作します。

そして配列と文字列もオブジェクトです。
